% LaTeX companion of 01-review-on-basic-concepts.typ. Both formats are editable.
\chapter{Review on Basic Concepts}

\section{Subspace and direct sum}

% qlnotes: kind=definition; id=subspace
\begin{definition}{\kn{Subspace}}\label{subspace}

vector space 的 subset \(U \subset V\) 为一个 subspace，if 它满足条件：

\begin{enumerate}
\tightlist
\item
  包含 0
\item
  对 addition 和 scalar multiplication 闭合
\end{enumerate}

\end{definition}

两个 subset 的和就是各取一个元素相加的所有情况.\\
很显然我们知道：

\begin{proposition}

两个 subspace \(U_{1},U_{2}\) 的 sum \(U_{1} + U_{2}\) 也是一个 subspace, 并且

\[\dim\left( {U_{1} + U_{2}} \right) \leq \dim\left( U_{1} \right) + \dim\left( U_{2} \right)\]

且 \(U_{1} + U_{2}\) 是同时包含 \(U_{1}\) 和 \(U_{2}\) 的 \(V\) 的最小 subspace.

\end{proposition}

显然可以随便和。同一个 \(U\) 自己和自己的和就是自己。所以 subspace sum 这个概念比较大，没什么用。我们需要用 direct sum 来作为一个小一点但是更有用的概念，表达出一种垂直的 subspace 的直观.

% qlnotes: kind=definition; id=direct sum
\begin{definition}{\kn{Direct sum}}\label{direct sum}

如果 \(U_{1} + U_{2} + \ldots + U_{m}\) 中的任意元素 \(v\)，都存在唯一的 \(v_{k} \in U_{k}\) for each \(k\) 使得 \(v = \sum_{k}v_{k}\)，就称 \(U_{1} + \ldots + U_{m} = \oplus_{i = 1}^{m}U_{i}\) 为一个 direct sum.

\end{definition}

我们显然发现：

\begin{proposition}

\[\dim\left( {\oplus_{i = 1}^{m}U_{i}} \right) = \sum\limits_{i = 1}^{m}\dim\left( U_{i} \right)\]

\end{proposition}

我们发现，其实可以 direct sum 的 subspaces 是 ``垂直的''，意思是:

\begin{theorem}

\(U_{1} + U_{2} + \ldots + U_{m}\) 是一个 direct sum (这几个空间``垂直'') iff 任取 \(u_{1},u_{2},\ldots,u_{m}\) 分别来自 \(U_{1},U_{2},\ldots,U_{m}\)，它们都 lin. ind.

\end{theorem}

并且：

\begin{theorem}

\(U_{1} + U_{2}\) 为一个 direct sum iff \(U_{1} \cap U_{2} = \left\{ 0 \right\}\).

\end{theorem}

\begin{qlremark}{Remark}

实际上两个 subspace 的交集里只要有一个非 0 点，那么这个点 span 的整个 dim 为 1 的线都在交集里.

\end{qlremark}

\begin{qlnote}{Note}

\(\mathbb{F}^{n} = \oplus_{i = 1}^{n}\text{span}\left( e_{i} \right)\)

\end{qlnote}
